\section{Mixing Console}

\subsection{Signal Summation}
See page 139 of teacher’s book.

Mixing console general pipeline: we have $N$ channels with signal processing blocks.  
At the end of the audio processing pipeline, we usually want to sum up different channels together:
\[
y(t) = \sum_{k=1}^{N} a_k \, x_k(t)
\]
\begin{figure}[H]
    \centering
    \includegraphics[width=0.9\linewidth]{Mixing_Console_Gen.png}
    \caption{General Overview of the signal summation}
\end{figure}

The weights $\{a_k\}_{k=1}^N$ are from the mixer of the channels, such that $a_k \in [0,1], \ \forall k \in [1,N]$.

To sum voltages, we can connect them in series. However, this is not satisfying if we have a lot of channels. Instead, we’ll put current generators in parallel.
\begin{figure}[H]
    \centering
    \includegraphics[width=0.8\linewidth]{Mixing_Console_limitation.png}
    \caption{Illustration of the limitations of summing voltages.}
\end{figure}

In an analog point of view, the voltages sum is obtained through the series of component voltages:
\[
\begin{cases}
v_0 = v_1 - iR_1 \\
v_0 = v_2 + iR_2
\end{cases}
\quad \Rightarrow \quad
i= \frac{v_1 - v_2}{R_1 + R_2}
\quad \Rightarrow \quad
v_0 = v_1 - \frac{v_1 - v_2}{R_1+R_2} \, R_1
\]
Thus:
\[
v_0 = v_1 \left( 1-\frac{R_1}{R_1+R_2}\right) + v_2\frac{R_1}{R_1+R_2} 
\quad \Rightarrow \quad v_0 = v_1 \frac{R_2}{R_1+R_2} + v_2\frac{R_1}{R_1+R_2}\neq v_1 + v_2
\]
As we do not obtain the desired result $v_0 = v_1 + v_2$, an idea is to sum currents instead of voltages, which requires a transconductance stage.
 
\begin{figure}[H]
    \centering
    % Left side: explanation + equations
    \begin{minipage}[c]{0.45 \linewidth}
    The simplest way is to create a current, as a function of a given voltage using a resistor, $i = f(v)$.
    
    NB: we don’t want the current to go to the ground, but we still want a ground reference. We talk about \textbf{virtual ground}. 

    Thus, we’ll use a transconductance circuit to generate a current proportional to the target voltage. To get rid of the ground issue, we’ll use an OPAMP. 
    
    The negative feedback configuration will solve the problem since $\varepsilon = 0$. This is understood as summing (negative) feedback amplifier.
    \end{minipage}\hfill
    \begin{minipage}[c]{0.5\linewidth}
        \centering
        \includegraphics[width=\linewidth]{Mixing_Console_Summing_Currents.png}
    \caption{Summing currents.}
    \end{minipage}
\end{figure}


With a single input, we obtain:  $ v_{o} = -v_1 \frac{R_f}{R_1} $.
\begin{figure}[H]
\centering
\includegraphics[width=0.3\linewidth]{Mixing_Console_1_input.png}
\caption{OPAMP circuit with 1 input.}
\end{figure}


\begin{figure}[H]
    \centering
    % Left side: explanation + equations
    \begin{minipage}[c]{0.65 \linewidth}
    With 2 inputs, we obtain: 
    \[
    v_{0}|_{v_2=0} = -v_1 \frac{R_f}{R_1}, 
    \quad v_{0}|_{v_1=0} = -v_2 \frac{R_f}{R_2}
    \]
    Thus we obtain that:
    \[
    v_0 = v_{0}|_{v_2=0} + v_{0}|_{v_1=0} = -v_1 \frac{R_f}{R_1} -v_2 \frac{R_f}{R_2} = -R_f \left( \frac{v_1}{R_1} + \frac{v_2}{R_2} \right)
    \]
    In this way, multiple currents can be summed together into the feedback branch, and superposition occurs:  
    \end{minipage}\hfill
    \begin{minipage}[c]{0.3\linewidth}
        \centering
        \includegraphics[width=\linewidth]{Mixing_Console_2_inputs.png}
        \caption{OPAMP circuit with 2 input.}
    \end{minipage}
\end{figure}

NB: we are not talking about mixing yet, that is a slightly different operation (with multipliers).  

A problem occurs when we consider too many channels entering the pin ``$-$'' of the OPAMP:  
\[
G_{LOOP} = A_0 \cdot \beta = A_0 \cdot \frac{R_{IN}}{N \cdot R_F} \propto \frac{1}{N} \quad \Rightarrow \quad \lim_{N \to +\infty} G_{LOOP} = 0
\]
Where the $\beta$ factor (voltage divider between input and feedback) is such that:
\[
\beta = \frac{R_{IN}}{N \cdot R_F}
\]
However the above configuration works only if $G_{LOOP} \gg 1$.

To solve this issue, we usually stack up different summing stages. For example:

\begin{figure}[H]
    \centering
    \includegraphics[width=0.5\linewidth]{Mixing_Console_Multiple_Summing_layers.png}
    \caption{Illustration of multiple 2-summing layers.}
\end{figure}



\subsection{Crosstalk}
Crosstalk is the coupling between a signal varying in time and others because of coupling capacitors $C_{STRAY}$.  
They are due to the close placement of metal elements (wire, plaque...).  
They create short circuits between nodes at high frequencies.

Thanks to the current summation of signals, we don’t get an effective coupling, leading to channel separation.

In the circuit below, L stands for Left channel, R for Right channel, C for Centre channel.

\begin{figure}[H]
    \centering
    \includegraphics[width=0.8\linewidth]{Mixing_Console_Crosstalk.png}
    \caption{Illustration of the Crosstalk.}
\end{figure}

\clearpage

\subsection{Panning potentiometer structure}

The potentiometers for panning are such that, when $R_{POT,L}$ goes up, $R_{POT,R}$ goes down and vice-versa.  

We usually introduce a voltage buffer (gain 1), characterized by an output resistor $R$.

\begin{figure}[H]
    \centering
    \includegraphics[width=0.2\linewidth]{Pan_Pot_1.png}
    \caption{Model of the panning potentiometer.}
\end{figure}

\begin{figure}[H]
    \centering
    % Left side: explanation + equations
    \begin{minipage}[c]{0.6 \linewidth}
    The total power will be given by $P_{TOT} = P_L + P_R$, and
    \[
    P = \frac{v^2}{R} \quad \Rightarrow \quad P_{TOT} = \frac{v_L^2}{R} + \frac{v_R^2}{R}
    \]
    \textbf{Case 1: Full-left}
    \[
    v_L = v_s, \quad v_R = 0 
    \quad \Rightarrow \quad
    P_{TOT} = \frac{v_s^2}{R}
    \]
    \textbf{Case 2: Full-right}
    \[
    v_L = 0, \quad v_R = v_s 
    \quad \Rightarrow \quad
    P_{TOT} = \frac{v_s^2}{R}
    \]
    \textbf{Case 3: Centre}
    \[
    v_L = v_R = \frac{v_s}{2}
    \quad \Rightarrow \quad
    P_{TOT} = \frac{v_L^2}{R} + \frac{v_R^2}{R}
    = \frac{v_s^2}{4R} + \frac{v_s^2}{4R}
    = \frac{v_s^2}{2R}
    \]
    \end{minipage}\hfill
    \begin{minipage}[c]{0.35\linewidth}
        \centering
        \includegraphics[width=\linewidth]{Pan_Pot_Circuit_1.png}
    \caption{Panning potentiometer circuit.}
    \end{minipage}\hfill
\end{figure}

At centre, the power is half compared to full-right or full-left.  
This creates a power gap, meaning a different amount of loudness is perceived.



\begin{figure}[H]
    \centering
    % Left side: explanation + equations
    \begin{minipage}[c]{0.5 \linewidth}
    To correct it, design of the panning is crucial. Ideally, we want:
    \[
    P_{TOT} = P_L + P_R = \frac{v_L^2}{R} + \frac{v_R^2}{R}
    = \frac{1}{R}\left(v_L^2 + v_R^2\right) = c^{st}
    \]
    NB: $ v_L = f_L(v_s), \quad v_R = f_R(v_s) $ which varies with the potentiometer position.
    
    Assuming a rotary potentiometer, we have:
    \[
    \begin{cases}
    v_L = v_s \cdot \sin(\theta) \\
    v_R = v_s \cdot \cos(\theta)
    \end{cases}
    \]
    \end{minipage}\hfill
    \begin{minipage}[c]{0.2\linewidth}
        \centering
        \includegraphics[width=\linewidth]{Pan_Pot_Bad.png}
    \caption{Basic yet inefficient potentiometer power graph.}
    \end{minipage}\hfill
    \begin{minipage}[c]{0.2\linewidth}
    \centering
    \includegraphics[width=\linewidth]{Pan_Pot_Good.png}
    \caption{Rectified panning potentiometer power graph.}
    \end{minipage}
\end{figure}
Thus, we obtain:
\[
P_{TOT} = \frac{1}{R} v_s^2 \left[ \sin^2(\theta) + \cos^2(\theta) \right]
= \frac{1}{R} v_s^2, \quad \forall \theta \in [0,2\pi]
\]
Thus, the problem here is solved using a rotary potentiometer.

We can also use compensated potentiometers: we add bypass resistors $R_b$ that give smoother transmission.
\begin{figure}[H]
    \centering
    \includegraphics[width=0.3\linewidth]{Pan_Pot_Bypass_Resistor.png}
    \caption{Scheme of the panning potentiometer circuit with 2 bypass resistors.}
\end{figure}

\subsection{Comments on the Mixing Console Presentation}

\begin{itemize}
    \item \textbf{+48V Phantom Power}: supplies phantom power to microphones that require it (typically condenser microphones).  
    It acts as a pure voltage shift, ensuring the mic receives the necessary bias for operation.

    \item \textbf{Transformers (two sets)}: remove any DC component from the signal path, allowing only the AC (audio) component to pass through.  
    They also provide galvanic isolation, helping reduce noise and interference.

    \item \textbf{Rng (Range) Button}: when engaged, this bypasses the first amplifier stage, routing the signal directly to the second amplifier.  
    \textit{Note}: the sensitivity scale on the knob must be read carefully, as two values are indicated (inner and outer markings), each corresponding to a different gain reference.

    \item \textbf{Phase Button (Ø symbol)}: inverts the signal’s polarity.  
    Useful when multiple microphones capture the same source but are subject to phase cancellation. By inverting one signal, destructive interference can be mitigated.  
    (\textit{See OneNote examples on phase inversion in mic setups.})

    \item \textbf{Band-Pass Filters (BP)}: placed at the end of the pre-amplification stage to attenuate unwanted low and high frequency components before equalization.

    \item \textbf{Post Pre-Amp Equalizer}:  
    \begin{itemize}
        \item Inner numbers on the knob’s crown: cutoff frequency of the \textbf{LPF} (Low-Pass Filter).
        \item Outer numbers on the knob’s crown: cutoff frequency of the \textbf{HPF} (High-Pass Filter).
    \end{itemize}
    This dual marking system allows precise control over both low- and high-frequency shaping.

    \item \textbf{PFL (Pre-Fader Listen)}: monitoring mode where the signal is taken \textit{before the fader}.  
    This allows the operator to check levels or tone without being affected by the fader’s position in the mix.

    \item \textbf{AFL (After-Fader Listen)}: monitoring mode where the signal is taken \textit{after the fader}.  
    This reflects exactly what is being sent to the mix bus, including any adjustments made with the fader.
\end{itemize}


\subsection{Notes from the slides}

A \textbf{mixing console} (or mixing desk) is the central element in both recording and live-sound systems.  
It allows multiple audio signals to be preamplified, equalized, combined, monitored, and routed toward various outputs.  
Its architecture exemplifies the integration of analog and digital electronics for high-fidelity and flexible sound processing.

Two principal families exist:
\begin{itemize}
  \item \textbf{Analog consoles:} perform all signal processing in the continuous voltage domain.
  \item \textbf{Digital consoles:} include analog-to-digital (A/D) and digital-to-analog (D/A) conversion, with internal DSP processing.
\end{itemize}


\begin{figure}[H]
    \centering
    \includegraphics[width=0.7\linewidth]{Analog_vs_Digital_Mixing_Console.png}
    \caption{Comparison between analog and digital console architectures.  
Analog consoles operate fully in the voltage domain, while digital ones use conversion and DSP for processing.}
\end{figure}

A professional mixing console typically comprises:
\begin{itemize}
  \item \textbf{Input modules:} preamplifiers, filters, equalizers, pan, mute, and fader.
  \item \textbf{Group and auxiliary buses:} allow submixing, effects routing, and monitor sends.
  \item \textbf{Output modules:} summing amplifiers, master faders, and output processing.
  \item \textbf{Control and automation:} VCAs (Voltage Controlled Amplifiers) or digital control layers.
\end{itemize}

Each input can be routed to one or several outputs through a programmable \textbf{mixing matrix}.

\begin{figure}[H]
    \centering
    % Left side: explanation + equations
    \begin{minipage}[c]{0.45 \linewidth}
    \centering
        \includegraphics[width=\linewidth]{Mixing_Console_Matrix.png}
    \caption{Simplified signal flow of a mixing console. Each channel can feed multiple outputs through routing and grouping.}
    \end{minipage}\hfill
    \begin{minipage}[c]{0.45\linewidth}
        \centering
        \includegraphics[width=\linewidth]{Mixing_Console_Routing_Matrix.png}
    \caption{Simplified signal flow of the mixing console's routing matrix.}
    \end{minipage}
\end{figure}


\paragraph{Input Module}
The input module is the fundamental building block. It contains:
\begin{enumerate}
  \item \textbf{Preamplifier:} provides gain (typically +20 to +60 dB) to bring microphone-level signals up to line level.
  \item \textbf{Filters:} high-pass and low-pass filters with adjustable cutoffs (e.g., 20–800\,Hz, 1–20\,kHz).
  \item \textbf{Equalizer (EQ):} usually a 4-band parametric EQ with variable frequency, gain, and quality factor $Q$.
  \item \textbf{Auxiliary sends:} direct signal copies to monitor or effects buses.
  \item \textbf{Pan control:} distributes the signal between left and right (L–R) outputs.
  \item \textbf{Fader:} adjusts the output signal level following a logarithmic law.
  \item \textbf{Mute and routing switches:} determine the channel’s active buses.
\end{enumerate}


\paragraph{Tone Shaping and Filtering}
After preamplification, the signal passes through:
\begin{itemize}
  \item A \textbf{high-pass filter} (HPF) to remove rumble and low-frequency noise.
  \item A \textbf{low-pass filter} (LPF) to reduce high-frequency interference.
  \item A \textbf{parametric equalizer} for tone control and correction.
\end{itemize}

Low-frequency (LF) and high-frequency (HF) sections often allow switching between \textbf{bell} and \textbf{shelf} responses:
\begin{itemize}
  \item \textbf{Bell response:} localized boost/cut centered around $f_0$.
  \item \textbf{Shelf response:} gradual boost/cut below or above a threshold frequency.
\end{itemize}

\begin{figure}
    \centering
    \includegraphics[width=0.5\linewidth]{Mixing_Console_Bell_vs_Shelf.png}
    \caption{Parametric equalizer responses: bell and shelf characteristics for LF and HF bands.}
\end{figure}


\paragraph{Auxiliary (AUX) Lines}
Auxiliary lines allow parallel routing for:
\begin{itemize}
  \item Effects sends (reverb, delay),
  \item Monitor mixes,
  \item Recording or broadcast feeds.
\end{itemize}

They can be configured as:
\begin{itemize}
  \item \textbf{Pre-fader} – independent from channel fader adjustments,
  \item \textbf{Post-fader} – follow channel fader level,
  \item \textbf{Pre-EQ or Post-EQ} – depending on whether they are affected by the equalizer section.
\end{itemize}

A professional console (e.g., Soundcraft Series V) may include up to 12 auxiliary buses, grouped in stereo or mono pairs.

\paragraph{Groups and Submixes}
Input channels can be assigned to one or several \textbf{groups}.  
A group collects and controls a set of related channels (e.g., drum mics, backing vocals) using a single fader.  
Groups can be panned between odd/even buses for stereo submixing. \textbf{Pan-switch behavior:}
\begin{itemize}
  \item When the pan switch is ON: groups receive the post-pan-pot signals (odd groups = left, even groups = right).
  \item When the pan switch is OFF: groups receive the pre-pan-pot signal (centered).
\end{itemize}

\paragraph{Monitoring: PFL and AFL Lines}
Monitoring systems use two specific routing modes:
\begin{itemize}
  \item \textbf{PFL (Pre-Fader Listen):} allows listening to the signal before the fader, unaffected by level changes.
  \item \textbf{AFL (After-Fader Listen):} monitors the post-fader signal, reflecting mix balance.
\end{itemize}

A \textbf{SOLO} switch can combine PFL and AFL functions, isolating one channel in the monitor output.

\paragraph{Voltage Controlled Amplifiers (VCA)}
In professional consoles, manual fader movements can be replaced or paralleled by a \textbf{Voltage Controlled Amplifier (VCA)} system.  
The signal gain is governed by a control voltage $V_C$:
\[
\frac{v_{out}}{v_{in}} = G(V_C),
\]
where $G(V_C)$ is the gain law defined by the control circuitry.

VCAs allow:
\begin{itemize}
  \item grouped control over several channels without altering signal summing,
  \item automation or remote fader control,
  \item accurate recall of mix levels.
\end{itemize}



\begin{figure}[H]
    \centering
    % Left side: explanation + equations
    \begin{minipage}[c]{0.65 \linewidth}
    \paragraph{Output Modules and Routing}
The output section aggregates all buses and provides final routing and level management:
\begin{itemize}
  \item \textbf{LCR (Left–Center–Right)} main outputs,
  \item \textbf{Group outputs} (e.g., 8),
  \item \textbf{Auxiliary outputs} (e.g., 12),
  \item \textbf{Matrix outputs} for complex distribution to external destinations.
\end{itemize}
\medskip
Each output may include:
\begin{itemize}
  \item send/return loops for external processors,
  \item master fader,
  \item signal monitoring and metering circuits.
\end{itemize}
    \end{minipage}\hfill
    \begin{minipage}[c]{0.3\linewidth}
        \centering
        \includegraphics[width=\linewidth]{Mixing_Console_Output_Matrix.png}
    \caption{Simplified schematic of a 16×10 matrix section.}
    \end{minipage}
\end{figure}


\paragraph{Master Section}
The master section combines multiple functions:
\begin{itemize}
  \item \textbf{L–R output module:} includes oscillator, talkback, and main bus routing.
  \item \textbf{Monitor and headphone section:} controls PFL/AFL monitoring and auxiliary returns.
  \item \textbf{Matrix and external routing:} sums groups and auxiliaries into custom mixes.
\end{itemize}

\paragraph{Inside the Console: Preamplifier and VCA Implementation}
The console’s preamplifier uses low-noise transistors (e.g., pnp BJT 2SB737) in differential configuration for excellent common-mode rejection.  
Two cascaded stages provide gain and headroom, while the following VCA stage (e.g., THAT2181) allows smooth gain control.

\paragraph{Analog vs. Digital Consoles}
\textbf{Analog consoles} rely on continuous electronic processing and are often valued for their sound character and tactile workflow.  
\textbf{Digital consoles}, instead, offer:
\begin{itemize}
  \item recallable settings and automation,
  \item integration with DAWs,
  \item onboard digital effects,
  \item flexible virtual routing.
\end{itemize}

Despite different implementations, both follow the same conceptual structure:  
\[
\text{\emph{input → processing → grouping → summing → output}
}
\]


